\subsection{A quartic for which the two tests differ}
Take $L$ to be the splitting field of $f(X)=X^4-X-1$. Its discriminant is $-283$. Its reduction modulo $2$ is irreducible: it has no linear factor, and its remainder modulo $X^2+X+1$ is $1$. Hence its Galois group contains a $4$-cycle. The cubic resolvent $Y^3+4Y-1$ has no rational root. Its three roots correspond to the partitions of the four quartic roots into pairs; this gives an action transitive on three points. Thus the group order is divisible by both $4$ and $3$. A subgroup of order $12$ in $S_4$ would be the index-two subgroup $A_4$, which has no $4$-cycle. It follows that the group is $S_4$. The known abelian field is $F=\mathbb Q(\sqrt{-283})$. The intermediate normal field $M=L^{V_4}$ is the splitting field of the resolvent.

For each $p\ne283$, factor $f$ modulo $p$. The one-sided criterion gives the following two distinct problems, with a separate centering constant for each row.
\begin{center}\small
\begin{tabular}{@{}llcl@{}}\toprule
Allowed factor degrees & First primes & Density & Zero problem seen\\\midrule
$(3,1)$ & $7,11,13,23,29,\ldots$ & $1/3$ & RH for $\zeta_M$\\
$(3,1)$ or $(2,2)$ & $7,11,13,23,29,\ldots$ & $11/24$ & RH for $\zeta_L$\\\bottomrule
\end{tabular}\end{center}
The first prime distinguishing the rows is $71$, of type $(2,2)$. Of the $167$ unramified primes up to $1000$, the first row selects $59$ and the second selects $82$. These counts follow by taking $\gcd(f,X^{p^j}-X)$ modulo $p$ for $j=1,2$. For the second row, the missing theorem is the lower estimate
\[
 \sum_{\substack{p\le x\\p\text{ in the second row}}}
       (\log p)(p^{-1}-x^{-1})
 \ \ge\ \frac {11}{24}\log x+c_L-C_L x^{-1/2}
 \quad\text{for all sufficiently large }x .
\]
Here $c_L$ is prescribed by the Laurent expansion, while $C_L$ may depend on the field. A geometric bound of this form would prove RH for this particular nonabelian field. The estimates in the pinned source are for the progression and quadratic-twist weights in its construction. A comparison of twisted lattices giving the displayed quartic weight is still missing.

